\documentclass[../main.tex]{subfiles}
\begin{document}

Chương~\ref{chapter:Experiment} đã trình bày chi tiết cách hệ thống được thiết kế và hiện
thực hoá. Chương này nhìn lại những giải pháp kỹ thuật mang tính \textbf{đóng góp nổi bật}
nhất --- đã được liệt kê tóm tắt ở mục~\ref{sec:dong-gop} --- và phân tích sâu hơn về lý do
chúng tạo nên sự khác biệt của hệ thống so với các cách làm thông thường, kèm theo những ưu
điểm và giới hạn còn tồn tại. Để tránh trùng lặp, phần mô tả cơ chế hoạt động chi tiết được
tham chiếu tới các mục tương ứng trong Chương~\ref{chapter:Experiment}.

\section{Cơ chế mã QR HMAC xoay theo thời gian}
\label{sec:dg-qr}

Phần lớn các hệ thống điểm danh bằng mã QR hiện nay sử dụng \textbf{mã QR tĩnh}: mỗi buổi học
(hoặc mỗi lớp) gắn với một mã cố định. Điểm yếu cố hữu của cách làm này là mã có thể bị chụp
màn hình và gửi cho người vắng mặt, khiến việc ``điểm danh hộ'' từ xa trở nên dễ dàng --- đây
chính là hình thức gian lận phổ biến nhất mà Chương~\ref{chapter:Related_works} đã chỉ ra.

Đóng góp của đồ án là thay thế mã tĩnh bằng \textbf{mã QR động được ký bằng HMAC-SHA256 và làm
mới sau mỗi 30 giây} (cơ chế và thuật toán sinh mã đã được trình bày ở
mục~\ref{sec:hmac-rotating} và Thuật toán~\ref{algo:qr}). Giá trị cốt lõi nằm ở ba tính chất
kết hợp: (i) mã chỉ có hiệu lực trong cửa sổ ngắn (90 giây) nên ảnh chụp nhanh chóng hết hạn;
(ii) chữ ký HMAC khiến không thể tự tạo ra một mã hợp lệ nếu không có khoá bí mật --- vốn được
cô lập ở phía máy chủ; và (iii) mốc thời gian gắn trong mã ngăn việc phát lại (replay).

Một điểm tinh tế nhưng quan trọng là cơ chế \textbf{đồng bộ đồng hồ giữa nhiều nguồn hiển thị}:
mã QR trên điện thoại giảng viên và mã trên màn hình máy chiếu (mục~\ref{sec:projector}) cùng
được sinh theo một mốc thời gian thống nhất, nhờ đó luôn khớp nhau ở mỗi chu kỳ làm mới. Cửa
sổ dung sai 90 giây (gấp ba lần chu kỳ) được lựa chọn để dung hoà sai lệch đồng hồ giữa các
thiết bị mà vẫn giữ thời gian sống của mã đủ ngắn. Nhờ cách tiếp cận này, hệ thống vô hiệu hoá
hình thức gian lận phổ biến nhất mà không đòi hỏi bất kỳ phần cứng chuyên dụng nào, chỉ dùng
chính chiếc điện thoại sẵn có của người dùng.

\section{Cổng máy chiếu và cơ chế pairing one-shot}
\label{sec:dg-projector}

Khi áp dụng mã QR động vào lớp học đông, một vấn đề thực tiễn nảy sinh: mã hiển thị trên màn
hình điện thoại của giảng viên quá nhỏ để sinh viên ở cuối phòng có thể quét, trong khi việc
chuyền tay điện thoại vừa bất tiện vừa mất an toàn. Đây là rào cản thường bị bỏ qua trong các
nghiên cứu nhưng lại quyết định khả năng sử dụng thực tế của hệ thống.

Đóng góp của đồ án là tính năng \textbf{Cổng máy chiếu} (mục~\ref{sec:projector}): một trang
web mở trên máy chiếu của lớp sẽ hiển thị mã QR động cỡ lớn, được liên kết với phiên điểm danh
của giảng viên thông qua một \textbf{mã ghép nối dùng một lần} (one-shot pairing token) dài
128 bit, có thời hạn 60 giây.

Điểm mới về mặt an toàn nằm ở \textbf{quy tắc chuyển trạng thái một chiều}: mã ghép nối chỉ đi
từ trạng thái \texttt{pending} sang \texttt{paired} đúng một lần, không thể quay lại cũng
không thể bị ghép nối lần thứ hai. Thiết kế này ngăn chặn nguy cơ một bên thứ ba chiếm quyền
hiển thị (hijack) màn hình máy chiếu để phát mã giả. Nhờ đó, tính năng vừa giải quyết bài toán
khả dụng cho lớp đông, vừa giữ nguyên các bảo đảm an toàn của cơ chế mã QR động --- một sự kết
hợp giữa trải nghiệm người dùng và bảo mật mà các ứng dụng điểm danh thông thường ít khi cung
cấp.

\section{Hệ thống chống gian lận nhiều lớp}
\label{sec:dg-fraud}

Hầu hết các hệ thống điểm danh hiện có chỉ dựa trên \textbf{một yếu tố xác minh duy nhất} (một
lần quét mã, hoặc một lần nhận diện khuôn mặt). Hệ quả là chỉ cần vượt qua yếu tố đó là gian
lận thành công --- một điểm thất bại đơn lẻ (single point of failure).

Đóng góp của đồ án là một quy trình điểm danh \textbf{đa yếu tố bốn bước} (mã QR, khuôn mặt,
xác minh ngang hàng và định vị) được tổng hợp bằng điểm tin cậy
(mục~\ref{sec:attendance-flow}), thiết kế theo nguyên lý \textbf{phòng thủ theo chiều sâu}
(defense in depth): mỗi lớp chống lại một hình thức gian lận khác nhau, và để qua mặt hệ thống,
người gian lận phải vượt qua đồng thời nhiều lớp. Sự tương ứng giữa từng lớp xác minh và hình
thức gian lận bị ngăn chặn được tóm tắt trong Bảng~\ref{tab:defense-layers}.

\begin{table}[H]
\centering
\small
\caption{Các lớp xác minh và hình thức gian lận bị ngăn chặn}
\label{tab:defense-layers}
\begin{tabular}{|p{5cm}|p{8.5cm}|}
\hline
\textbf{Lớp xác minh} & \textbf{Hình thức gian lận bị ngăn chặn} \\
\hline
Mã QR động (HMAC) & Chụp và chia sẻ lại mã để điểm danh hộ \\ \hline
Xác minh khuôn mặt & Nhờ người khác điểm danh thay \\ \hline
Xác minh ngang hàng & Điểm danh từ xa khi không thực sự có mặt \\ \hline
Định vị GPS & Điểm danh khi đang ở ngoài khu vực lớp học \\ \hline
Phân tích hành vi phía máy chủ & Thông đồng có hệ thống (cùng nhóm, xác minh quá nhanh\ldots) \\ \hline
\end{tabular}
\end{table}

Trong các lớp trên, \textbf{xác minh ngang hàng} (peer-to-peer) là đóng góp giàu tính mới mẻ
nhất. Thay vì chỉ tin vào một bằng chứng do chính sinh viên cung cấp, hệ thống yêu cầu sinh
viên trao đổi mã QR với tối thiểu ba bạn xung quanh; việc xác minh chéo này tạo ra một
\textbf{bằng chứng về sự hiện diện vật lý} mà một yếu tố đơn lẻ (mã hay khuôn mặt) khó tạo ra
được --- rất khó để ``được ba người trong phòng xác minh'' nếu không thực sự có mặt. Bổ trợ
cho lớp phòng ngừa này, một hàm phân tích chạy phía máy chủ rà soát các mẫu hành vi đáng ngờ
(Bảng~\ref{tab:fraud}) như luôn xác minh với cùng một nhóm hay tốc độ xác minh bất thường, qua
đó bổ sung thêm một lớp \textbf{phát hiện} bên cạnh các lớp \textbf{ngăn chặn}.

Cần thẳng thắn nhìn nhận giới hạn: xác minh ngang hàng vẫn có thể bị một nhóm sinh viên thông
đồng cùng có mặt để xác minh khống cho nhau. Tuy nhiên, chính lớp phân tích hành vi được thiết
kế để phát hiện kiểu thông đồng lặp lại này, và việc kết hợp nhiều lớp khiến chi phí gian lận
tăng cao hơn đáng kể so với hệ thống chỉ dùng một yếu tố.

\section{Kiến trúc ba vai trò và phân quyền}
\label{sec:dg-role}

Một thách thức của đề tài là phải thực thi phân quyền chặt chẽ trên một nền tảng bị ràng buộc:
Zalo Mini App không có máy chủ xác thực truyền thống, còn hệ thống vận hành phần lớn trên gói
dịch vụ miễn phí của Firebase. Trong bối cảnh đó, việc bảo đảm mỗi vai trò chỉ làm đúng phần
việc của mình là không hiển nhiên.

Đóng góp của đồ án là một \textbf{kiến trúc ba vai trò} (sinh viên, giảng viên, quản trị) với
phân quyền được thực thi ngay tại tầng cơ sở dữ liệu bằng quy tắc bảo mật của Firestore, kết
hợp cơ chế cấp token tùy biến (mục~\ref{sec:auth-design}, Bảng~\ref{tab:phanquyen}). Nguyên
tắc thiết kế cốt lõi là \textbf{ngăn tự nâng quyền} (privilege escalation): tài khoản tạo từ
phía thiết bị chỉ được mang vai trò \texttt{student}, còn vai trò giảng viên hay quản trị chỉ
có thể được gán bởi hàm phía máy chủ. Đồng thời, mọi nghiệp vụ nhạy cảm --- ký mã QR, tính
điểm tin cậy, phát hiện gian lận --- đều được tập trung xử lý ở máy chủ để không thể bị can
thiệp từ thiết bị người dùng.

Điểm đáng chú ý là kiến trúc này đạt được mô hình \textbf{đặc quyền tối thiểu} (least
privilege) và chống leo thang quyền mà \textbf{không cần một máy chủ ứng dụng riêng}, qua đó
phù hợp với ràng buộc và chi phí của nền tảng. Giới hạn hiện tại --- một số quy tắc còn dựa
trên cấu trúc dữ liệu thay vì danh tính đầy đủ do hạn chế của gói miễn phí --- đã được nhận
diện và có thể khắc phục khi nâng cấp gói dịch vụ, như đã đề cập ở mục~\ref{sec:bao-mat}.

\section{Tuân thủ chính sách nền tảng Zalo}
\label{sec:dg-zalo}

Nhiều sản phẩm thử nghiệm trong môi trường học thuật thường bỏ qua các yêu cầu về chính sách
nền tảng và quyền riêng tư, dẫn tới việc không thể thực sự phát hành tới người dùng. Đồ án này
xác định \textbf{khả năng triển khai thực tế} là một mục tiêu, nên việc tuân thủ chính sách
trở thành một đóng góp có ý nghĩa thực tiễn.

Cụ thể, ứng dụng được xây dựng để tuân thủ chính sách của Zalo Mini App: không yêu cầu cấp
quyền tại thời điểm mở ứng dụng mà chỉ xin quyền khi người dùng chủ động thao tác, và xin theo
từng phạm vi riêng biệt (thông tin cá nhân, số điện thoại). Hệ thống cung cấp đầy đủ Điều
khoản sử dụng và Chính sách quyền riêng tư, đồng thời xử lý \textbf{dữ liệu sinh trắc học một
cách có trách nhiệm}: khuôn mặt chỉ được lưu dưới dạng véc-tơ đặc trưng 128 chiều chứ không lưu
ảnh gốc, trên cơ sở có sự đồng thuận của người dùng.

Hướng tới khả năng áp dụng cho nhiều cơ sở đào tạo khác nhau, sản phẩm còn được điều chỉnh
theo hướng trung lập về thương hiệu. Những yếu tố này tuy không nằm ở phần ``lõi thuật toán''
nhưng lại là điều kiện cần để một hệ thống điểm danh có thể đi từ nguyên mẫu tới một sản phẩm
dùng được trong thực tế.

\subsection*{Kết chương}

Chương này đã phân tích năm đóng góp nổi bật của đồ án: cơ chế mã QR động HMAC có đồng bộ đa
thiết bị, tính năng Cổng máy chiếu với ghép nối một chiều an toàn, hệ thống chống gian lận
nhiều lớp với điểm nhấn là xác minh ngang hàng, kiến trúc phân quyền chặt chẽ trên nền tảng bị
ràng buộc, và việc tuân thủ chính sách nền tảng hướng tới triển khai thực tế. Xuyên suốt các
đóng góp này là một tư tưởng nhất quán: nâng cao tính trung thực của điểm danh bằng cách kết
hợp nhiều lớp bảo vệ chi phí thấp, trong khi vẫn giữ trải nghiệm đơn giản cho người dùng và
khả năng triển khai thực tế. Phần kết luận sau đây sẽ nhìn lại toàn bộ kết quả và bàn về các
định hướng mở rộng tiếp theo.

\end{document}
